\chapter{Introduction}
\label{ch:intro}

\section{What hdl-vet is for}

A generator writes your SystemVerilog. You need a few changes in its output:
a new port on the top module, a tie-off that is no longer a constant, a
technology cell in place of a generic module. If you edit the generated files
by hand, you lose the edits the next time the generator runs.

hdl-vet keeps the edits in \term{scripts}. A script is a short Python file. When the
generator runs again, you run the scripts again, and you get the same edits on
the new files. Each edit is checked. If the result is not what the script
asked for, hdl-vet writes nothing.

hdl-vet uses \texttt{pyslang}, the Python binding of the slang compiler, to
understand the design: its modules, instances, ports, parameters and
connections. It makes each edit as a replacement of a range of bytes in the
source text. After the edits, it loads the edited design again and compares it
with the design you started from. This last step is the \term{sign-off}.

\section{The words used in this manual}

Each word has one meaning in this manual.

\begin{description}
\item[input tree] The folder that holds the design as the generator wrote it. It is
  the deepest folder that contains every source file, every include folder and the
  \emph{base} of the file list (the folder that relative file names in the file list start
  from; chapter~\ref{ch:first}). hdl-vet never changes the input tree. The file list
  itself may lie outside the input tree.
\item[output tree] A new folder that hdl-vet writes. It has the same layout as the
  input tree. It holds the files with your edits, and the files without edits as
  exact copies. A file list that lies outside the input tree is written at the top of
  the output tree.
\item[file and path] A \emph{file} is a file on disk. A \emph{path} is a place in the
  design: a chain of instance names that starts with the top module, for example
  \code{soc_top.u_periph.u_timer0}. This manual uses \emph{file} for the first and
  \emph{path} only for the second.
\item[script] A Python file that defines the function \code{apply(vet)}. It names the
  edits you want.
\item[operation] One edit that hdl-vet knows, for example \apiname{Module.add_port}
  or \apiname{Instance.connect}. A script is a list of calls to operations.
\item[applied] An operation is \term{applied} when it changes the text.
\item[skipped] An operation is \term{skipped} when the design already shows the
  result it asks for. A skipped operation changes nothing. This is why you can
  run a script a second time.
\item[drift] The design is in a state that the operation does not expect: it
  is neither the state before the edit nor the state after it. The operation
  stops with an error. Chapter~\ref{ch:scripts} explains how \code{was=} detects drift.
\item[queued] An operation does not change the text at once. It records the edit in a
  \term{queue}. hdl-vet applies the queue to the text when it next loads the design.
\item[baseline] The design as the compiler first loaded it, before any edit. The
  sign-off compares the final design with the baseline.
\item[elaborated, copy] The compiler builds the design from the top module
  downwards. A module is \term{elaborated} once for each place where it is
  instantiated, and each of these is a \term{copy} of the module. \code{timer} is
  instantiated twice in the example design, so it has two copies. A module that
  nothing instantiates below the top module is not elaborated.
\item[session] The state of a run: the text of every file with the edits so far, the
  queue, the expectations, the checks, the log, the notes, the added files and the
  scripts that ran. You can save a session and continue it later
  (chapter~\ref{ch:session}).
\item[checked operation] An operation that tries its edit on the design before it
  returns, and is refused if the result is wrong (chapter~\ref{ch:scripts}).
\item[expectation] A statement that a structural change must happen
  (\code{expect}) or may happen (\code{permit}) during the run. Operations
  register their expectations themselves. You write one by hand only when you
  edit the text yourself.
\item[sign-off] The final check of a run. hdl-vet loads the edited design and
  compares it with the baseline. If anything is wrong, it writes no output tree.
  Chapter~\ref{ch:signoff} lists everything that it checks.
\end{description}

\section{What you can rely on}

\begin{itemize}
\item The input tree is never changed.
\item The output tree is a copy of the input tree with your edits.
\item Text that no operation edits stays byte-identical. hdl-vet does not
  reformat anything. An operation inserts new text with the line ending of the
  file it edits.
\item Every file that changes gets one extra comment line,
  \code{// Modified by hdl-vet; scripts: ...}, with the names and the hashes
  of the scripts that ran.
\item The sign-off decides whether the output tree is written. If it fails,
  the exit status is 1 and no output tree is written.
\item The same input and the same scripts give the same output, byte for byte.
\end{itemize}

\section{What hdl-vet does not do}

\begin{itemize}
\item It does not check what your logic does. The sign-off compares the
  structure of the design (instances, ports, parameters, connections) and the
  diagnostics of the compiler. It does not simulate. A change of logic that
  you make with \apiname{Vet.splice} passes the sign-off unless you declare it.
  Chapter~\ref{ch:limits} gives an example.
\item It does not edit instances in generate blocks, instance arrays, or
  instances that come from a macro.
\item It sees one configuration of the design: the defines of your file list.
\item It does not write the design from nothing. It edits what a generator
  wrote.
\end{itemize}

Chapter~\ref{ch:limits} collects what is missing and what is known to be wrong.
No other chapter hints at the future.

\section{How to read this manual}

Chapter~\ref{ch:first} takes you through a first project, from start to end.
Read it first. Chapters~\ref{ch:scripts} to~\ref{ch:signoff} explain scripts,
operations and the sign-off. Chapters~\ref{ch:cli} to~\ref{ch:session} describe
the command line, the project file and saved sessions. Chapter~\ref{ch:python}
lists the functions you use when you call hdl-vet from your own Python code.

Every example in this manual is a file, and a test runs it. There are two kinds.
A \term{script} is shown with its file name. A \term{transcript} shows commands
that you type, and what they print. In a transcript, a line that starts with
\texttt{\$ } is a command. The lines below it are its output. A last line
\code{[exit status 2]} shows a non-zero exit status.
A transcript runs in a copy of the folder \file{examples}, which holds the
example design. Where a file name of your own machine would appear, the transcript
shows \code{<examples>}. The program prints its own version as
\code{<version>}.

The example design is small: \code{soc_top} instantiates a reset synchroniser
\code{rst_sync} and a block \code{periph_wrap}, which instantiates the counter
\code{timer} twice. Chapter~\ref{ch:first} shows its files.
